% GENERATED FILE. Edit canonical lessons, then run npm run generate:books.
% canonical-source-hash: fnv1a64:69f37d1997fb95b0
% canonical-lessons: ZH-C517-tianliang, ZH-C517-guang-light, ZH-W517-guang, ZH-C517-faguang, ZH-C517-xingguang, ZH-C517-huanghun

\chapter{Light in the Dark}
\label{ch:zh-light-in-the-dark}
\hlchaptermodality{517}

\begin{chapteropening}
\textbf{By the end of this chapter:} \emph{I can talk about light, things that shine, starlight, dusk and the time it gets light.}

Complete the last lesson of chapter 517: I can talk about light, things that shine, starlight, dusk and the time it gets light.
\end{chapteropening}

\section[tiānliàng]{tiānliàng — to get light at dawn}

\label{lesson:ZH-C517-tianliang}



\begin{quote}
Before the new one: say the Mandarin for the woods, then the Mandarin for to ripen.
\end{quote}

\subsection*{You'll want to know: tiānliàng}
\emph{tiānliàng} — \textquotedblleft{}to get light at dawn\textquotedblright{}.

\begin{quote}
\emph{wàipó shuō dōngtiān qī diǎn cái tiānliàng} — Grandma says it doesn't get light until seven in winter.
\end{quote}

Its characters come later; for now it is a word for the ear.

\begin{cousinweb}
\emph{tiānliàng} is built from \zh{天} 'the sky' and \emph{liàng} 'bright', both inherited from Middle Chinese, \emph{then} and \emph{ljangH} (Guangyun, in Baxter's transcription), which glosses the second 'bright' and 'clear'.
\end{cousinweb}

\begin{grammarlens}[title={What you've built}]
One more word for saying how it was.
\end{grammarlens}

\subsection*{Your turn}
\begin{itemize}
\raggedright
  \item \emph{Say it:} \emph{tiānliàng}
  \item \emph{Say it:} \emph{tiānliàng}, once more
  \item \emph{Say it:} \emph{wàipó shuō dōngtiān qī diǎn cái tiānliàng} — Grandma says it doesn't get light until seven in winter.
  \item \emph{Recall:} say the Mandarin for the woods, then the Mandarin for to ripen, then say \emph{tiānliàng} again
\end{itemize}

\begin{tcolorbox}[breakable,colback=teal!4,colframe=teal!35!black,title={Before you move on}]
What does tiānliàng mean? (\textquotedblleft{}To get light at dawn\textquotedblright{}.) Say it once more.
\end{tcolorbox}
\section[\texorpdfstring{\zh{光} (guāng)}{光 (guāng)}]{\zh{光} (guāng) — light from the sun or a lamp}

\label{lesson:ZH-C517-guang-light}



\begin{quote}
Before the new one: say the Mandarin for to ripen, then the Mandarin for to get light at dawn.
\end{quote}

\subsection*{You'll want to know: \zh{光}}
\textbf{\zh{光}} — \emph{guāng} — \textquotedblleft{}light from the sun or a lamp\textquotedblright{}.

\begin{quote}
\textbf{\zh{停电以后房间里没有光了}} \emph{tíngdiàn yǐhòu fángjiān lǐ méiyǒu guāng le} — After the power went off, there was no light in the room.
\end{quote}

The new character, in pieces: \zh{光} — six strokes in three pieces: on top a short vertical with a dot on each side, three strokes; a long bar across the middle; and \zh{儿} at the foot, two strokes. The next lesson puts it in your hand.

\begin{cousinweb}
\textbf{\zh{光}} (\emph{guāng}) is inherited from Middle Chinese \emph{kwang} (Guangyun, in Baxter's transcription), glossed 'bright'; the dictionary also lists a departing-tone reading, \emph{kwangH}, for adding colour. The noun 'light' is not glossed there.
\end{cousinweb}

\begin{grammarlens}[title={What you've built}]
One more word for saying how it was.
\end{grammarlens}

\subsection*{Your turn}
\begin{itemize}
\raggedright
  \item \emph{Say it:} \emph{guāng}
  \item \emph{Say it:} \emph{guāng}, once more
  \item \emph{Say it:} \emph{tíngdiàn yǐhòu fángjiān lǐ méiyǒu guāng le} — After the power went off, there was no light in the room.
  \item \emph{Recall:} say the Mandarin for to ripen, then the Mandarin for to get light at dawn, then say \emph{guāng} again
\end{itemize}

\begin{tcolorbox}[breakable,colback=teal!4,colframe=teal!35!black,title={Before you move on}]
What does \zh{光} mean? (\textquotedblleft{}Light from the sun or a lamp\textquotedblright{}.) Say it once more.
\end{tcolorbox}
\section[\texorpdfstring{\zh{光} (guāng)}{光 (guāng)}]{\zh{光} — three short strokes, a bar, then \zh{儿}}

\label{lesson:ZH-W517-guang}



\begin{quote}
The word first, then the character inside it.

\begin{itemize}
\raggedright
  \item \emph{Recall:} say \emph{guāng} — \textbf{R1}, one lesson back, before it decays
\end{itemize}

\begin{quote}
\zh{儿} \zh{先}
\end{quote}

\emph{Read it:} each one aloud, with nothing said first — \emph{ér}, \emph{xiān}
\end{quote}

\subsection*{Script — three short strokes, a bar, then \zh{儿}}
\hlblockfigure{\detokenize{figures/ZH-W517-guang-filmstrip.pdf}}{How \zh{光} is written, stroke by stroke}

\begin{quote}
\zh{光}
\end{quote}

\textbf{Three pieces, top to bottom.} Count the pieces, not the strokes.

\textbf{The top piece, three strokes}: a short vertical down the middle, a dot on its left, then a short stroke on its right falling to the left.

\textbf{One long bar} across the middle, under them.

\textbf{The bottom piece, \zh{儿}, two strokes}, as you wrote it: a stroke falling to the left, then a vertical that bends right along the foot and hooks up.

3 + 1 + 2 = six.

\subsection*{Your turn}
\begin{itemize}
\raggedright
  \item \emph{Write it:} \zh{光}, all six, with the model in front of you
  \item \emph{Cover:} write it again from memory, and count the pieces, not the strokes
\end{itemize}

\begin{tcolorbox}[breakable,colback=teal!4,colframe=teal!35!black,title={Before you move on}]
How many strokes does \zh{光} have? (\textbf{Six}.) What sits at its foot? (\textbf{\zh{儿}, two strokes}.)

Source: \href{https://raw.githubusercontent.com/chanind/hanzi-writer-data/68d10a4b21150cae5e1ebbd223eed289cf32d90c/data/\%E5\%85\%89.json}{Hanzi Writer Data \zh{光}, pinned stroke paths and medians}.
\end{tcolorbox}
\section[\texorpdfstring{\zh{发光} (fāguāng)}{发光 (fāguāng)}]{\zh{发光} (fāguāng) — to give off light; to shine}

\label{lesson:ZH-C517-faguang}



\begin{quote}
Before the new one: say the Mandarin for to get light at dawn, then the Mandarin for light from the sun.
\end{quote}

\subsection*{You'll want to know: \zh{发光}}
\textbf{\zh{发光}} — \emph{fāguāng} — \textquotedblleft{}to give off light; to shine\textquotedblright{}.

\begin{quote}
\textbf{\zh{星星为什么会发光}} \emph{xīngxing wèishénme huì fāguāng} — Why do the stars shine?
\end{quote}

\begin{cousinweb}
\textbf{\zh{发光}} (\emph{fāguāng}) is built from \zh{发} 'to send out' and \zh{光} 'light', both inherited from Middle Chinese, \emph{pjot} and \emph{kwang} (Guangyun, in Baxter's transcription), which glosses the first 'to start' and 'to open out' and the second 'bright'.
\end{cousinweb}

\begin{grammarlens}[title={What you've built}]
One more word for saying how it was.
\end{grammarlens}

\subsection*{Your turn}
\begin{itemize}
\raggedright
  \item \emph{Say it:} \emph{fāguāng}
  \item \emph{Say it:} \emph{fāguāng}, once more
  \item \emph{Say it:} \emph{xīngxing wèishénme huì fāguāng} — Why do the stars shine?
  \item \emph{Recall:} say the Mandarin for to get light at dawn, then the Mandarin for light from the sun, then say \emph{fāguāng} again
\end{itemize}

\begin{tcolorbox}[breakable,colback=teal!4,colframe=teal!35!black,title={Before you move on}]
What does \zh{发光} mean? (\textquotedblleft{}To give off light; to shine\textquotedblright{}.) Say it once more.
\end{tcolorbox}
\section[\texorpdfstring{\zh{星光} (xīngguāng)}{星光 (xīngguāng)}]{\zh{星光} (xīngguāng) — starlight}

\label{lesson:ZH-C517-xingguang}



\begin{quote}
Before the new one: say the Mandarin for light from the sun, then the Mandarin for to give off light.
\end{quote}

\subsection*{You'll want to know: \zh{星光}}
\textbf{\zh{星光}} — \emph{xīngguāng} — \textquotedblleft{}starlight\textquotedblright{}.

\begin{quote}
\textbf{\zh{山里的星光最好看}} \emph{shān lǐ de xīngguāng zuì hǎokàn} — The starlight is loveliest up in the hills.
\end{quote}

\begin{cousinweb}
\textbf{\zh{星光}} (\emph{xīngguāng}) is built from \zh{星} 'a star' and \zh{光} 'light', both inherited from Middle Chinese, \emph{seng} and \emph{kwang} (Guangyun, in Baxter's transcription), which glosses the second 'bright'.
\end{cousinweb}

\begin{grammarlens}[title={What you've built}]
One more word for saying how it was.
\end{grammarlens}

\subsection*{Your turn}
\begin{itemize}
\raggedright
  \item \emph{Say it:} \emph{xīngguāng}
  \item \emph{Say it:} \emph{xīngguāng}, once more
  \item \emph{Say it:} \emph{shān lǐ de xīngguāng zuì hǎokàn} — The starlight is loveliest up in the hills.
  \item \emph{Recall:} say the Mandarin for light from the sun, then the Mandarin for to give off light, then say \emph{xīngguāng} again
\end{itemize}

\begin{tcolorbox}[breakable,colback=teal!4,colframe=teal!35!black,title={Before you move on}]
What does \zh{星光} mean? (\textquotedblleft{}Starlight\textquotedblright{}.) Say it once more.
\end{tcolorbox}
\section[huánghūn]{huánghūn — twilight; the fading light of evening}

\label{lesson:ZH-C517-huanghun}



\begin{quote}
Before the new one: say the Mandarin for to give off light, then the Mandarin for starlight.
\end{quote}

\subsection*{You'll want to know: huánghūn}
\emph{huánghūn} — \textquotedblleft{}twilight; the fading light of evening\textquotedblright{}.

\begin{quote}
\emph{wǒ zuì xǐhuan huánghūn de shíhou qù sànbù} — Twilight is my favourite time for a walk.
\end{quote}

Its characters come later; for now it is a word for the ear.

\begin{cousinweb}
\emph{huánghūn} is built from \emph{huáng} 'yellow' and \emph{hūn}, both inherited from Middle Chinese, \emph{hwang} and \emph{xwon} (Guangyun, in Baxter's transcription), which glosses the first as the colour of the centre and, after an older dictionary, the second as the sun gone dark.
\end{cousinweb}

\begin{grammarlens}[title={What you've built}]
One more word for saying how it was.
\end{grammarlens}

\subsection*{Your turn}
\begin{itemize}
\raggedright
  \item \emph{Say it:} \emph{huánghūn}
  \item \emph{Say it:} \emph{huánghūn}, once more
  \item \emph{Say it:} \emph{wǒ zuì xǐhuan huánghūn de shíhou qù sànbù} — Twilight is my favourite time for a walk.
  \item \emph{Recall:} say the Mandarin for to give off light, then the Mandarin for starlight, then say \emph{huánghūn} again
\end{itemize}

\begin{tcolorbox}[breakable,colback=teal!4,colframe=teal!35!black,title={Before you move on}]
What does huánghūn mean? (\textquotedblleft{}Twilight; the fading light of evening\textquotedblright{}.) Say it once more.
\end{tcolorbox}
